\section{Idea 2: Conservative Value Learning}

\begin{frame}{Problem: Offline Values Can Be Overconfident}
  \Large
  To compute useful weights, we need value estimates.

  \vspace{0.35cm}
  \normalsize
  But value estimation is harder offline:
  \begin{itemize}
    \item online learning can test actions and correct bad guesses,
    \item offline learning cannot collect new feedback,
    \item unsupported actions may receive overly optimistic predicted values.
  \end{itemize}

  \vspace{0.35cm}
  \begin{block}{Question}
    How can we learn values without trusting predictions where the data is weak?
  \end{block}
  \note{This is the only place to briefly reintroduce why offline value learning is hard. Keep it short.}
\end{frame}

\begin{frame}{Central Idea: Be Pessimistic Where Evidence Is Weak}
  \centering
  \begin{tikzpicture}[>=Latex,x=1cm,y=0.85cm]
    \draw[->,thick] (0,0) -- (10,0) node[right] {action};
    \foreach \x/\h in {1.4/1.0,2.2/1.7,3.0/1.3,3.9/2.0,4.8/1.5} {
      \draw[fill=Panel] (\x-0.12,0) rectangle (\x+0.12,\h);
    }
    \node[above] at (3.1,2.25) {evidence};
    \draw[SignalGreen,ultra thick] (1.0,1.9) -- (5.2,1.9);
    \node[right] at (5.25,1.9) {trusted};
    \draw[WarmGold,dashed,ultra thick] (7.6,0) -- (7.6,2.8);
    \node[above] at (7.6,2.8) {unsupported};
    \draw[WarmGold,ultra thick] (7.0,0.45) -- (8.2,0.45);
    \node[right] at (8.25,0.45) {pessimistic};
  \end{tikzpicture}

  \vspace{0.35cm}
  \Large
  Conservative value learning lowers confidence away from the logged data.

  \vspace{0.2cm}
  \normalsize
  The policy should not be attracted to actions whose value is high only because the model is guessing.
  \note{This slide conveys the CQL principle visually. Pessimism is a guardrail against unsupported optimism.}
\end{frame}

\begin{frame}{Representative Method: CQL}
  \Large
  Conservative Q-Learning changes the value-learning objective.

  \vspace{0.35cm}
  \normalsize
  Its principle is simple:
  \begin{itemize}
    \item fit observed transitions in the dataset,
    \item penalize high values assigned to actions not well supported by the dataset,
    \item prefer policies whose improvement is backed by evidence.
  \end{itemize}

  \vspace{0.35cm}
  \begin{block}{What to remember}
    Do not trust unsupported actions.
  \end{block}

  \vspace{0.1cm}
  \small
  Representative algorithm: Conservative Q-Learning (CQL).
  \note{Avoid the CQL objective. The audience should remember the pessimism principle, not the optimization details.}
\end{frame}

\begin{frame}{Remaining Problem: Can We Avoid the Risk Entirely?}
  \Large
  CQL makes unsupported actions less attractive.

  \vspace{0.35cm}
  \normalsize
  But it still starts from the same danger:

  \vspace{0.2cm}
  \[
    \text{searching over actions can expose value errors.}
  \]

  \vspace{0.35cm}
  \begin{block}{Next question}
    Can we estimate values without maximizing over unsupported actions at all?
  \end{block}
  \note{This transition motivates IQL. The next idea is not just more pessimism; it avoids the problematic maximization step.}
\end{frame}
